% Document CEM1. Canonical proof manuscript; companion outline CEO1.
\documentclass[11pt,letterpaper,reqno]{amsart}
\usepackage[bibliography]{manuscript-style}
\newcommand{\R}{\mathbb R}
\newcommand{\Z}{\mathbb Z}
\newcommand{\N}{\mathbb N}
\newcommand{\E}{\mathbb E}
\newcommand{\dd}{\mathrm d}


\newcommand{\C}{\mathbb C}

\title[Gaussian domination and continuum limits]
 {Gaussian domination and structural properties of continuum limits}
\date{}
\hypersetup{pdftitle={Gaussian domination and structural properties of continuum limits},pdfauthor={}}
\begin{document}
\begin{abstract}
Gaussian domination is preserved under weak convergence of laws of
tempered fields. Under a common $L^2$ Gaussian bound, the generating
functions of every finite family of Schwartz pairings converge
locally uniformly, as do their logarithms near the origin.
Consequently all joint moments and cumulants converge. We also
prove a lattice compactness theorem under the weaker assumption of
a uniform second-moment bound. Cell averaging and Hermite coefficient
estimates give subsequential continuum laws. We prove that approximate
linear identities, geometric symmetries and reflection inequalities
pass to these limits, and establish tensor Schwinger-function bounds.
For chamber QED3, the resulting laws are closed, translation and
cubically invariant, and reflection positive at the nine lattice
mirror directions at every height. We also recover the real lattice
law associated with a prescribed fixed-spacing volume limit.
\end{abstract}
\maketitle
\tableofcontents

\section{Introduction}\label{sec:introduction}
Let $\mathcal S=\mathcal S(\R^d;\R^r)$ be the real Schwartz space
and $\mathcal S'$ its continuous dual. Write $e_\alpha$ for the
standard basis of $\R^r$. For the normalized real Hermite functions
$h_k$, $k\in\N_0^d$, set
\[
 (1-\Delta+|y|^2)h_k=\lambda_k h_k,\qquad
 \lambda_k=2|k|_1+d+1,\qquad |k|_1=k_1+\cdots+k_d.
\]
For $t>0$ we use the separable Hilbert space
\begin{equation}\label{eq:hermite-space}
 H_{-t}=\{F\in\mathcal S':\|F\|_{-t}<\infty\},\qquad
 \|F\|_{-t}^2=
 \sum_{\alpha=1}^r\sum_{k\in\N_0^d}
       \lambda_k^{-2t}|F(h_ke_\alpha)|^2.
\end{equation}
All weak convergence of probability laws on this space refers to
its norm topology.

Extend each real distribution complex linearly to
$\mathcal S_{\mathbb C}=\mathcal S(\R^d;\mathbb C^r)$.
For a law $P$, its smeared $n$-point correlation is
$S_n^P(f_1,\ldots,f_n)=\E_P\prod_{j=1}^n F(f_j)$ whenever
this expectation exists.

\begin{theorem}[Gaussian domination and convergence of correlations]
\label{thm:gaussian-convergence}
Fix integers $d,r\ge1$ and a real number $t>0$. Let $P_j$ and $P$ be Borel probability
laws on $H_{-t}$ such that $P_j\Longrightarrow P$. Suppose
\begin{equation}\label{eq:sequence-gaussian}
 \E_{P_j}e^{F(f)}\le e^{\|f\|_2^2/2}
 \qquad(j\ge1,\ f\in\mathcal S).
\end{equation}
Then
\begin{equation}\label{eq:limit-continuum-G}
 \E_P e^{F(f)}\le e^{\|f\|_2^2/2}
 \qquad(f\in\mathcal S).
\end{equation}
For every finite family $f_1,\ldots,f_m\in\mathcal S_{\mathbb C}$,
the functions
\[
 Z_j(z)=\E_{P_j}\exp\left(\sum_{a=1}^m z_aF(f_a)\right),
 \qquad
 Z(z)=\E_P\exp\left(\sum_{a=1}^m z_aF(f_a)\right)
\]
are entire on $\mathbb C^m$. They satisfy $Z_j\to Z$ locally
uniformly, with every complex derivative. In particular, for
every $n\ge1$ and $f_1,\ldots,f_n\in\mathcal S_{\mathbb C}$,
\begin{equation}\label{eq:correlation-convergence}
 S_n^{P_j}(f_1,\ldots,f_n)\longrightarrow
 S_n^P(f_1,\ldots,f_n).
\end{equation}
On the common domain
\begin{equation}\label{eq:zero-free-domain}
 D_f=\left\{z\in\mathbb C^m:
                 \sum_{a=1}^m|z_a|\|f_a\|_2<2/5\right\},
\end{equation}
all $Z_j$ and $Z$ satisfy $|Z_j-1|<1/2$ and $|Z-1|<1/2$.
Their principal logarithms converge locally uniformly on $D_f$,
with every derivative. Thus all joint cumulants converge, where
\[
 \kappa_n^P(f_1,\ldots,f_n)
 =\left.\partial_{z_1}\cdots\partial_{z_n}
       \log\E_P\exp\left(\sum_{a=1}^n z_aF(f_a)\right)
   \right|_{z=0},
\]
and $\kappa_n^{P_j}$ is defined in the same way.
\end{theorem}

The exponential estimate controls the unbounded observables in
\eqref{eq:correlation-convergence}. More precisely, it gives
a uniform second-moment bound for each exponential in $Z_j$,
so weak convergence passes these expectations to the limit.
Local bounds for the entire generating functions then pass all
their derivatives. The bound $|Z_j-1|<1/2$ on
$D_f$ permits the same argument for their
logarithms. We give the proof in \cref{sec:gaussian}, following
the random-linear-form estimates and convergence argument in
\cite[Sections~5.1--5.2 and~6.2]{QED3Manuscript}.

The weakly convergent subsequences needed for this result can be
obtained from second moments alone. We state the lattice form
separately, since it does not require Gaussian domination.

\begin{theorem}[Continuum limits of lattice fields]\label{thm:lattice-continuum}
Fix integers $d,r\ge1$, offsets $t_1,\ldots,t_r\in[0,1]^d$,
and a set $I\subset(0,1]$ with $0$ in its closure. For each
$a\in I$, let $\nu_a$ be a Borel probability law on
$\R^{\Z^d\times\{1,\ldots,r\}}$ with the coordinate product
topology. Suppose there is a constant $0\le C<\infty$ such that
\begin{equation}\label{eq:lattice-second}
 \E_{\nu_a}\left|\sum_{\alpha,x}b_\alpha(x)u_\alpha(x)\right|^2
 \le C\sum_{\alpha,x}|b_\alpha(x)|^2
\end{equation}
for every $a\in I$ and every finitely supported real coefficient
array $b$. Define
\begin{equation}\label{eq:step-definition}
 (F_a)_\alpha(y)=a^{-d/2}u_\alpha(x)
 \quad\text{for }y\in Q_{a,\alpha,x}
 :=a\bigl(x+t_\alpha+[-1/2,1/2)^d\bigr).
\end{equation}
Then $F_a$ is almost surely a real tempered distribution, and
\begin{equation}\label{eq:smeared-second}
 \E_{\nu_a}|F_a(f)|^2\le C\|f\|_{L^2(\R^d;\R^r)}^2
 \qquad(f\in\mathcal S).
\end{equation}
For every fixed $t>d/2$, these fields have Borel realizations in
$H_{-t}$ whose laws are uniformly tight. Consequently, for every
sequence $a_j\in I$ tending to zero there are a subsequence
$a_{j_\ell}$ and a Borel probability measure $\mu$ on $H_{-t}$
such that
\[
 \mathcal L(F_{a_{j_\ell}})\ \Longrightarrow\ \mu.
\]
Along this subsequence the pairings with every finite family of
Schwartz tests converge jointly in distribution to their pairings
under $\mu$.
\end{theorem}

The hypothesis bounds the raw second moment of every finite linear
combination. Bounds on the individual coordinates suffice to realize
the step field, but the full inequality \eqref{eq:lattice-second}
is needed for \eqref{eq:smeared-second}. No compatibility between
the laws at different spacings is assumed. The conclusion is
subsequential existence; it does not assert uniqueness or that a
limit is nonzero.

There are two points in the proof. First, the cell coefficients
\begin{equation}\label{eq:cell-coefficients}
 b_{a,\alpha}(x)=a^{-d/2}\int_{Q_{a,\alpha,x}}f_\alpha(y)\,\dd y
\end{equation}
satisfy $\|b_a\|_{\ell^2}\le\|f\|_2$. Applying the lattice
bound to finite truncations therefore controls the smeared field
uniformly in spacing. In \cref{prop:step-realization}, a spatially
weighted lattice norm identifies these limits, for all tests at
once, with the samplewise step distribution. Second, the test estimate gives
\[
 \E\|F_a\|_{-s}^2\le Cr\sum_k\lambda_k^{-2s}<\infty
 \qquad(s>d/2).
\]
For $t>s$, a closed bounded ball in $H_{-s}$ is compact in $H_{-t}$:
the extra factor $\lambda_k^{-2(t-s)}$ makes its coefficient tails
uniformly small. Markov's inequality and Prokhorov's theorem then
give the asserted subsequences, as proved in \cref{prop:compactness}.
Choosing $d/2<s<t$ gives every exponent claimed in the theorem;
its proof concludes \cref{sec:compactness}.

This proof uses the Hermite compactness argument for curvature laws
in \cite[Sections~5.4 and~6.2]{QED3Manuscript}, together with the
weighted lattice realization in
\cite[Sections~5.3 and~6.1]{LatticeLimits}. We isolate the
second-moment assumptions used by those arguments. The dimension
and the number of components enter only through the lattice weight,
cell volume and Hermite multiplicities. The Hermite description of
Schwartz space and its dual is recalled from
\cite[Sections~2, 4 and~5]{RadhaThangavelu}; the probability compactness
result is Prokhorov's theorem in the form stated in
\cite[Theorem~1.1]{Sorensen}. The application in \cref{sec:qed}
uses the already constructed real-curvature laws of chamber QED3.
The next theorem supplies the identities, symmetries and reflection
forms retained by these limits.

For the structural assertion, identify the component space $\R^r$ with a real inner-product space \(V\), and let \(W\) be another
finite-dimensional real inner-product space.
Fix an orthogonal representation \(\rho:O(d)\to O(V)\).
For \(g(x)=Qx+b\), use the actions
\begin{equation}\label{eq:geometric-action}
 (g\cdot f)(x)=\rho(Q)f(Q^{-1}(x-b)),\qquad
 (g\cdot F)(f)=F(g^{-1}\cdot f).
\end{equation}
Invariance concerns the law as a random tempered distribution,
with the sigma-algebra generated by evaluations.
For an open set \(O\), write \(\Sigma_O\) for the sigma-algebra
generated by \(F(f)\), \(f\in C_c^\infty(O;V)\) real.
For an open half-space \(H^+\) and an \(r\)-invariant law \(P\),
where \(r\) reflects in its
boundary, reflection positivity means
\[
 \E_P[\overline{B(r\cdot F)}B(F)]\ge0
 \qquad\bigl(B\in L^2(\Sigma_{H^+},P)\bigr).
\]

\begin{theorem}[Identities, symmetries and reflection positivity]
\label{thm:structural-limit}
Let \(P_j\Rightarrow P\) weakly on \(H_{-t}\), and suppose
\begin{equation}\label{eq:common-gaussian}
 \E_{P_j}e^{F(f)}\le e^{\|f\|_2^2/2}
 \qquad(j\ge1,\ f\in\mathcal S\text{ real}).
\end{equation}
Let \(D^*:C_c^\infty(\R^d;W)\to\mathcal S\) be continuous and
linear. Let \(\mathcal G\) be a group of Euclidean isometries
containing all translations and a finite set \(\mathcal R\) of
reflections through the origin. Choose an open side \(H_r^+\) for
each \(r\in\mathcal R\). Assume:
\begin{enumerate}[label=(\roman*),leftmargin=*]
\item\label{hyp:identity}
For each \(\xi\in C_c^\infty(\R^d;W)\) there are real
\(v_j(\xi)\in\mathcal S\) such that \(F(v_j(\xi))=0\),
\(P_j\)-almost surely, and
\(\|v_j(\xi)-D^*\xi\|_2\to0\).
\item\label{hyp:symmetry}
For each \(g\in\mathcal G\) there are real linear maps
\(T_{j,g}:\mathcal S\to\mathcal S\) satisfying
\[
 \E_{P_j}e^{iF(T_{j,g}f)}=\E_{P_j}e^{iF(f)},\qquad
 \|T_{j,g}f-g\cdot f\|_2\longrightarrow0
 \quad(f\in\mathcal S).
\]
\item\label{hyp:reflection}
For every \(r\in\mathcal R\), every finite family of real
\(f_a\in C_c^\infty(H_r^+;V)\), and every \(z_a\in\C\),
\begin{equation}\label{eq:rp-approx}
 \sum_{a,b}\overline z_a z_b
 \E_{P_j}e^{iF(f_b-r\cdot f_a)}\ge0
\end{equation}
for all sufficiently large \(j\), depending on the family.
\end{enumerate}
Then \(P\) satisfies \(\E_Pe^{F(f)}\le e^{\|f\|_2^2/2}\)
for every real \(f\in\mathcal S\), is
\(\mathcal G\)-invariant, and almost every \(F\) satisfies
\(F(D^*\xi)=0\) simultaneously for all \(\xi\).
The law is reflection invariant and reflection positive at every
translate of every plane in \(\mathcal R\).
Along the same sequence, all joint moments and cumulants of every
fixed finite family of Schwartz evaluations converge to those of \(P\).
\end{theorem}

A linear identity is retained because its error vanishes in mean
square. Symmetry is expressed by characteristic functions, whose
change is bounded by the \(L^2\) error in the transformed test.
Reflection positivity passes through a finite sum of bounded
continuous characters; reflection invariance then extends the form
to the positive-side \(L^2\) space. These are the structural limit
arguments of \cite[Sections~6.3 and~6.5]{QED3Manuscript}, stated
here with their test-approximation hypotheses explicit. Their proofs
are given in \cref{sec:stability}.

Gaussian stability and convergence of moments and cumulants are
provided by \cref{thm:gaussian-convergence}.
The structural arguments require only
\(\E|F(f)|^2\le C\|f\|_2^2\), except that passage of the reflection
form requires no moment estimate. The common exponential hypothesis
supplies the quadratic control and the correlation conclusions.

The application in \cref{sec:qed} starts from specified
infinite-lattice laws. Their existence, reflection inequalities
and estimates are supplied by \cite[Theorem~A]{ChamberLattice}.
The fixed-spacing identification with prescribed volume limits is
proved in \cref{sec:recovery}, following
\cite[Section~6.4]{QED3Manuscript}.
The resulting symmetry group is the affine cubic group; the limit
argument introduces no additional rotations.


\section{Gaussian bounds and convergence of correlations}\label{sec:gaussian}
We first collect the estimates for one law. Their proofs are
the random-linear-form arguments of
\cite[Proposition~5.1 and Section~5.2]{QED3Manuscript}.
They apply in every finite dimension and for any finite number
of components, since they use only the $L^2$ norm of the tests.

\begin{proposition}\label{prop:subgaussian}
Let $Q$ be a probability law of real tempered distributions
with measurable test pairings, and suppose
\begin{equation}\label{eq:probability-G}
 \E_Q e^{F(f)}\le e^{\|f\|_2^2/2}\qquad(f\in\mathcal S).
\end{equation}
Then, for real $f$,
\begin{equation}\label{eq:field-variance}
 \E_Q F(f)=0,\qquad \E_Q|F(f)|^2\le\|f\|_2^2.
\end{equation}
For $p>0$ set $c_p=2\,2^{p/2}\Gamma(1+p/2)$. One has
\[
 \E_Q|F(f)|^p\le c_p\|f\|_2^p\qquad(f\in\mathcal S).
\]
For complex tests $f_1,\ldots,f_n$,
\begin{equation}\label{eq:joint-moment-bound}
 \E_Q\prod_{a=1}^n|F(f_a)|
       \le 2^{n/2}c_n\prod_{a=1}^n\|f_a\|_2.
\end{equation}
The generating function of each finite family of complex tests
is entire. For every $g\in\mathcal S_{\mathbb C}$,
\begin{align}
 \left|\E_Q e^{F(g)}\right|
       &\le e^{\|\operatorname{Re}g\|_2^2/2},
                      \label{eq:entire-generating-bound}\\
 \left|\E_Q e^{F(g)}-1\right|
       &\le e^{2\|g\|_2^2}-1. \label{eq:zero-free-bound}
\end{align}
\end{proposition}

\begin{proof}
For real $f$, put $Y=F(f)$ and $v=\|f\|_2$. Applying
\eqref{eq:probability-G} to $tf$ gives
\begin{align*}
 \E_Q e^{tY}&\le e^{t^2v^2/2} &&(t\in\R),\\
 \E_Q e^{t|Y|}&\le2e^{t^2v^2/2} &&(t\ge0).
\end{align*}
The second bound justifies differentiation at zero. Comparing
the linear terms in the first bound for both signs of $t$
gives $\E_QY=0$; comparing the quadratic terms gives
$\E_QY^2\le v^2$. If $v=0$, this also gives $Y=0$ almost
surely. If $v>0$, Markov's inequality with $t=u/v^2$,
applied to $Y$ and $-Y$, gives
\[
 Q(|Y|\ge u)\le2e^{-u^2/(2v^2)}\qquad(u>0).
\]
Integration yields
\[
 \E_Q|Y|^p
 =p\int_0^\infty u^{p-1}Q(|Y|>u)\,\dd u
 \le2\,2^{p/2}\Gamma(1+p/2)v^p.
\]
For $g=g_1+ig_2$ with real $g_1,g_2$ and integer $n\ge1$,
the triangle inequality in $L^n(Q)$ gives
\[
 \|F(g)\|_{L^n(Q)}
 \le c_n^{1/n}(\|g_1\|_2+\|g_2\|_2)
 \le\sqrt2\,c_n^{1/n}\|g\|_2.
\]
H\"older's inequality proves \eqref{eq:joint-moment-bound}.

On a compact polydisc in $\mathbb C^m$, the integrand
$\exp(\sum_a z_aF(f_a))$ is bounded in absolute value by
\[
 \exp\left(C\sum_a
       \bigl(|F(\operatorname{Re}f_a)|
             +|F(\operatorname{Im}f_a)|\bigr)\right)
\]
for a constant $C$ depending only on that polydisc. H\"older's
inequality and the exponential bound above make this majorant
integrable. The same bound on a larger polydisc justifies
differentiation under the expectation, so the generating
function is entire. Taking the modulus of its integrand and
applying \eqref{eq:probability-G} proves
\eqref{eq:entire-generating-bound}.

For the last estimate, put $\varphi(u)=e^u-1-u$. Since
$\E_Q F(g)=0$ and $|e^z-1-z|\le\varphi(|z|)$,
\[
 \left|\E_Q e^{F(g)}-1\right|
 \le\E_Q\varphi(|F(g)|).
\]
For $u\ge0$, $\varphi(u)\le2(\cosh u-1)$. Writing
$g=g_1+ig_2$, convexity and monotonicity on $[0,\infty)$
therefore give
\begin{align*}
 \E_Q\varphi(|F(g)|)
 &\le\tfrac12\E_Q\varphi(2|F(g_1)|)
       +\tfrac12\E_Q\varphi(2|F(g_2)|)\\
 &\le e^{2\|g_1\|_2^2}-1+e^{2\|g_2\|_2^2}-1\\
 &\le e^{2\|g\|_2^2}-1.
\end{align*}
This proves \eqref{eq:zero-free-bound}.
\end{proof}

\begin{proof}[Proof of \Cref{thm:gaussian-convergence}]
Schwartz evaluation is continuous on $H_{-t}$:
$|F(f)|\le\|F\|_{-t}\|f\|_t$, with the test norm
defined in the proof of \cref{prop:compactness}.
For real $f$, the function $F\mapsto\min\{R,e^{F(f)}\}$
is bounded and continuous. Weak convergence gives
\[
 \E_P\min\{R,e^{F(f)}\}
 =\lim_{j\to\infty}\E_{P_j}\min\{R,e^{F(f)}\}
 \le e^{\|f\|_2^2/2}.
\]
Letting $R\to\infty$ proves \eqref{eq:limit-continuum-G}.
Thus \cref{prop:subgaussian} applies to $P$ as well as to
every $P_j$. In particular, the generating functions in
the theorem are entire.

Fix $f_1,\ldots,f_m\in\mathcal S_{\mathbb C}$ and write
$g_z=\sum_a z_af_a$. The map
$F\mapsto\exp(F(g_z))$ is continuous, and
\begin{equation}\label{eq:exponential-ui}
 \E_{P_j}\bigl|e^{F(g_z)}\bigr|^2
 =\E_{P_j}e^{2F(\operatorname{Re}g_z)}
 \le e^{2\|\operatorname{Re}g_z\|_2^2}.
\end{equation}
The same inequality holds under $P$. For each $z$, it gives
uniform integrability of these exponentials. Truncating their
real and imaginary parts, passing the bounded continuous
expectations to the weak limit, and removing the truncations
therefore proves $Z_j(z)\to Z(z)$.

The right side of \eqref{eq:entire-generating-bound}, applied
to $g_z$, is bounded on each compact subset of $\mathbb C^m$,
uniformly in $j$. Cauchy's estimates on a larger polydisc give
uniform bounds for the first derivatives on a smaller one.
The functions $Z_j$ are consequently equicontinuous on compact
sets. A finite covering by small balls, together with pointwise
convergence and continuity of $Z$, upgrades the convergence to
local uniform convergence. Cauchy's formula on surrounding
polydiscs then gives local uniform convergence of every complex
derivative. Differentiation under the expectations identifies
the derivatives at zero with joint moments, proving
\eqref{eq:correlation-convergence}.

For $z\in D_f$, the triangle inequality gives
$\|g_z\|_2\le\sum_a|z_a|\|f_a\|_2<2/5$. Hence
\eqref{eq:zero-free-bound} gives
\[
 |Z_j(z)-1|,\ |Z(z)-1|
 \le e^{2\|g_z\|_2^2}-1
 < e^{8/25}-1<\tfrac12.
\]
All the images lie in the disk of radius $1/2$ about $1$,
on which the principal logarithm is analytic. Local uniform
convergence of $Z_j$ gives local uniform convergence of their
logarithms on $D_f$. Cauchy's formula again gives convergence
of every derivative, and in particular of every joint cumulant.
\end{proof}

\section{Realization of the step field}\label{sec:realization}
We first work at one fixed spacing. The estimate that makes the
step field a distribution may depend on this spacing. The estimate
on its second moments will retain the constant in
\eqref{eq:lattice-second}.

\begin{proposition}\label{prop:step-realization}
Fix $d,r\ge1$, $a>0$ and offsets $t_\alpha\in[0,1]^d$.
Suppose a Borel probability law $\nu$ on real lattice arrays
satisfies \eqref{eq:lattice-second} with constant $C$.
On one measurable set of full measure, the step function
\eqref{eq:step-definition} is tempered and, for every $f\in\mathcal S$,
\begin{equation}\label{eq:absolute-pairing}
 F_a(f)=\sum_{\alpha,x}b_{a,\alpha}(x)u_\alpha(x),\qquad
 \sum_{\alpha,x}|b_{a,\alpha}(x)u_\alpha(x)|<\infty.
\end{equation}
After setting $F_a=0$ off this set, every test pairing is measurable.
The finite truncations of \eqref{eq:absolute-pairing} converge in
$L^2(\nu)$, and
\[
 \E_\nu|F_a(f)|^2\le C\|b_a\|_{\ell^2}^2\le C\|f\|_2^2.
\]
\end{proposition}

\begin{proof}
Choose $p>d/2$ and put
\[
 \|u\|_{\ell^2_{-p}}^2
 =\sum_{\alpha,x}(1+|x|^2)^{-p}|u_\alpha(x)|^2.
\]
Apply \eqref{eq:lattice-second} to each coordinate. Tonelli's
theorem gives
\begin{equation}\label{eq:lattice-weighted}
 \E_\nu\|u\|_{\ell^2_{-p}}^2
 \le Cr\sum_{x\in\Z^d}(1+|x|^2)^{-p}<\infty.
\end{equation}
Thus $\Omega_p:=\{u:\|u\|_{\ell^2_{-p}}<\infty\}$ is a
measurable set of full measure, independent of the test function.

For $u\in\Omega_p$, weighted Cauchy--Schwarz yields
\begin{equation}\label{eq:weighted-pairing}
 \sum_{\alpha,x}|b_{a,\alpha}(x)u_\alpha(x)|
 \le\|u\|_{\ell^2_{-p}}
 \left(\sum_{\alpha,x}(1+|x|^2)^p
                       |b_{a,\alpha}(x)|^2\right)^{1/2}.
\end{equation}
The second factor is controlled by a Schwartz seminorm. Indeed,
Cauchy--Schwarz on each cube bounds its square by
\[
 \sum_{\alpha,x}(1+|x|^2)^p
             \int_{Q_{a,\alpha,x}}|f_\alpha(y)|^2\,\dd y
 \le C_{a,p,d}\int_{\R^d}(1+|y|^2)^p|f(y)|^2\,\dd y.
\]
Here $|x|\le a^{-1}|y|+\tfrac32\sqrt d$ on the cube, and
the cubes partition $\R^d$ for each component. The constant may
depend on $a,p,d$, but is uniform in the allowed offsets.
The right side is a continuous Schwartz seminorm squared.
Consequently \eqref{eq:absolute-pairing} defines a continuous
linear functional for every $u\in\Omega_p$. The same estimate
with $|f_\alpha|$ in the cell integrals shows absolute
integrability against the step function, so the functional is
precisely its distributional pairing. Finite coordinate sums are
measurable. Their pointwise limits on $\Omega_p$, extended by zero
outside it, prove measurability of the pairings.

For the uniform estimate, use Cauchy--Schwarz without the spatial
weight. Since each cube has volume $a^d$,
\begin{equation}\label{eq:cell-contraction}
 \sum_{\alpha,x}\left|a^{-d/2}
               \int_{Q_{a,\alpha,x}}f_\alpha(y)\,\dd y\right|^2
 \le\sum_{\alpha,x}\int_{Q_{a,\alpha,x}}|f_\alpha(y)|^2\,\dd y
 =\|f\|_2^2.
\end{equation}
Let $b_a^{(m)}$ be the restriction of $b_a$ to $|x|\le m$,
with every component retained. These restrictions converge to $b_a$
in $\ell^2$. Applying \eqref{eq:lattice-second} to
$b_a^{(m)}-b_a^{(n)}$ shows that their pairings form an
$L^2(\nu)$ Cauchy sequence. Their almost sure limit on $\Omega_p$
is \eqref{eq:absolute-pairing}, so their $L^2$ limit agrees with
$F_a(f)$. Passing to the limit in the finite quadratic bound and
then using \eqref{eq:cell-contraction} proves the assertion.
\end{proof}

\section{Compactness from Hermite coefficients}\label{sec:compactness}
The preceding proposition reduces the problem to a family of
random distributions with a common $L^2$ test estimate. We state
the resulting compactness criterion independently of the lattice.
The proof retains the coefficient argument of
\cite[Section~5.4]{QED3Manuscript} and
\cite[Section~6.1]{LatticeLimits}.

\begin{proposition}\label{prop:compactness}
Let $(X_i)_{i\in J}$ be a family of random real tempered
distributions in $\mathcal S'(\R^d;\R^r)$, defined on possibly
different probability spaces. Suppose $X_i(f)$ is
measurable for each $f\in\mathcal S$ and
\begin{equation}\label{eq:distribution-second}
 \E|X_i(f)|^2\le C\|f\|_2^2
 \qquad(i\in J,\ f\in\mathcal S)
\end{equation}
for one $0\le C<\infty$. For every $s>d/2$, the fields have
Borel realizations in $H_{-s}$ and
\begin{equation}\label{eq:hermite-expectation}
 \sup_i\E\|X_i\|_{-s}^2\le
 B_s:=Cr\sum_{k\in\N_0^d}\lambda_k^{-2s}<\infty.
\end{equation}
For every $t>s$ and $R>0$, the set
\[
 K_R=\{X\in H_{-s}:\|X\|_{-s}\le R\}
\]
is compact in $H_{-t}$ and
$\inf_i\mathbb P(X_i\in K_R)\ge1-B_s/R^2$.
The laws are uniformly tight in $H_{-t}$ and every sequence has
a weakly convergent subsequence there. The inclusion
$H_{-t}\hookrightarrow\mathcal S'$ is continuous for the strong
distribution topology.
\end{proposition}

\begin{proof}
We first recall the identification of the Hilbert space. If
$f_{\alpha k}=\langle f_\alpha,h_k\rangle_{L^2}$, the Hermite
expansion of $f\in\mathcal S$ converges in $\mathcal S$. The norm
$\|f\|_s$, defined by
\[
 \|f\|_s^2:=\sum_{\alpha,k}\lambda_k^{2s}|f_{\alpha k}|^2
\]
is a continuous Schwartz seminorm. Indeed, for real $s\ge0$ and
an integer $m\ge s$,
\[
 \|f\|_s\le\|(1-\Delta+|y|^2)^m f\|_2.
\]
These facts are the Hermite description
of Schwartz space in \cite[Sections~2, 4 and~5]{RadhaThangavelu}.
For any coefficients $c_{\alpha k}$ with
$\sum\lambda_k^{-2s}|c_{\alpha k}|^2<\infty$, the formula
\[
 f\longmapsto\sum_{\alpha,k}c_{\alpha k}f_{\alpha k}
\]
defines a tempered distribution by weighted Cauchy--Schwarz.
It has the prescribed Hermite coefficients. Conversely, the
Schwartz expansion identifies every $X\in H_{-s}$ with this
formula for $c_{\alpha k}=X(h_ke_\alpha)$. Thus $H_{-s}$ is
isometric to a weighted $\ell^2$ space and is separable and complete.

\emph{Expected norm and measurable realization.}
The bound \eqref{eq:distribution-second} gives
$\E|X_i(h_ke_\alpha)|^2\le C$. Tonelli's theorem therefore
gives \eqref{eq:hermite-expectation}. To check summability,
group the indices by $|k|_1=n$:
\begin{equation}\label{eq:hermite-trace}
 \sum_{k\in\N_0^d}\lambda_k^{-2s}
 =\sum_{n=0}^\infty
    \binom{n+d-1}{d-1}(2n+d+1)^{-2s}<\infty
 \qquad(2s>d).
\end{equation}
The multiplicity is of order $(1+n)^{d-1}$. On the measurable
full-measure event where the norm is finite, the finite Hermite
sums of $X_i$ converge in $H_{-s}$. Each sum is a Borel random
variable, since it has finitely many measurable coefficients.
Their limit, extended by zero off that event, is a Borel
$H_{-s}$-valued random variable. The coefficient identification
above shows that it agrees there with the given distribution
$X_i$, for every Schwartz test at once.

\emph{Compact balls.}
For $X\in K_R$, the tail of the weaker norm satisfies
\begin{equation}\label{eq:compact-tail}
 \sum_{\alpha,k:\lambda_k>M}\lambda_k^{-2t}
                    |X(h_ke_\alpha)|^2
 \le M^{-2(t-s)}R^2.
\end{equation}
Take a sequence $(X_n)$ in $K_R$. Each coefficient is bounded
by $R\lambda_k^s$, so a diagonal subsequence makes every
coefficient converge to a number $c_{\alpha k}$. Fatou's lemma
gives $\sum\lambda_k^{-2s}|c_{\alpha k}|^2\le R^2$.
These coefficients therefore define an $X\in K_R$. On the
finitely many indices with $\lambda_k\le M$, coefficient
convergence gives convergence in $H_{-t}$. On the remaining
indices, the squared norm of $X_n-X$ is at most
$4R^2M^{-2(t-s)}$, by \eqref{eq:compact-tail} applied to the
difference. First let $n\to\infty$ and then $M\to\infty$.
The chosen subsequence converges to $X$ in $H_{-t}$. This proves
compactness of $K_R$ in that metric space.

Markov's inequality and \eqref{eq:hermite-expectation} give
$\mathbb P(X_i\notin K_R)\le B_s/R^2$, uniformly in $i$.
Thus the laws are tight in the separable complete Hilbert space
$H_{-t}$. Prokhorov's theorem
\cite[Theorem~1.1]{Sorensen} gives their relative weak sequential
compactness.

Finally, for a bounded subset $D\subset\mathcal S$, weighted
Cauchy--Schwarz gives
\begin{equation}\label{eq:strong-inclusion}
 \sup_{f\in D}|X(f)|\le\|X\|_{-t}\sup_{f\in D}\|f\|_t.
\end{equation}
The last supremum is finite. These are the seminorms defining
the strong topology of $\mathcal S'$, so the inclusion is
continuous.
\end{proof}

\begin{proof}[Proof of \Cref{thm:lattice-continuum}]
Apply \cref{prop:step-realization} to each $\nu_a$. It gives
one tempered step distribution at that spacing, measurable test
pairings, and \eqref{eq:smeared-second} with the common constant
$C$. Fix $t>d/2$ and choose $d/2<s<t$.
\Cref{prop:compactness} gives Borel laws on $H_{-t}$,
uniform tightness and a weakly convergent subsequence along
any prescribed sequence $a_j\to0$. For any finite collection
$f_1,\ldots,f_m\in\mathcal S$, the map
$X\mapsto(X(f_1),\ldots,X(f_m))$ is continuous on $H_{-t}$
by \eqref{eq:strong-inclusion}. Composing bounded continuous
functions on $\R^m$ with this map proves the asserted joint
convergence. The continuous inclusion also regards the limit
as a Borel law on the space of tempered distributions.
\end{proof}

\section{The structural limit passages}\label{sec:stability}

\subsection{Linear identities and approximate symmetries}
Suppose \(P_j\Rightarrow P\) on \(H_{-t}\) and
\begin{equation}\label{eq:quadratic}
 \E_{P_j}|F(f)|^2\le C\|f\|_2^2
 \qquad(f\in\mathcal S\text{ real})
\end{equation}
with one finite \(C\).
Under \cref{hyp:identity}, exact annihilation of \(v_j(\xi)\) gives
\begin{equation}\label{eq:identity-residual}
 \E_{P_j}|F(D^*\xi)|^2
 \le C\|D^*\xi-v_j(\xi)\|_2^2\longrightarrow0.
\end{equation}
For \(M>0\), the function
\(\min\{|F(D^*\xi)|^2,M\}\) is bounded and continuous on \(H_{-t}\).
Its limiting expectation is zero, so \(F(D^*\xi)=0\),
\(P\)-almost surely. Choose a countable dense family of smooth tests
on each member of a compact exhaustion of \(\R^d\), and intersect
the corresponding full-measure events. Continuity of \(D^*\) and
of each sample distribution extends the identity to every compact
smooth test.

For symmetries, put \(\chi_j(f)=\E_{P_j}e^{iF(f)}\).
The inequality \(|e^{iu}-e^{iv}|\le|u-v|\) and
\eqref{eq:quadratic} imply
\begin{equation}\label{eq:character-continuity}
 |\chi_j(f)-\chi_j(h)|\le\sqrt C\,\|f-h\|_2
 \qquad(f,h\in\mathcal S\text{ real}).
\end{equation}
For a fixed \(g\) in \cref{hyp:symmetry}, this gives
\[
 |\chi_j(g\cdot f)-\chi_j(f)|
 \le\sqrt C\,\|g\cdot f-T_{j,g}f\|_2\longrightarrow0.
\]
Weak convergence passes both fixed-test characteristic functions
to the limit. Applying the equality to all real linear combinations
of any finite family proves equality of its joint laws.
These evaluations generate the sigma-algebra, so \(P\) is
\(g\)-invariant.

In particular, suppose \(P_j\) is invariant under translations in
\(a_j\Z^d\), \(a_j\to0\). Given \(b\), choose \(b_j\in a_j\Z^d\)
with \(b_j\to b\). The maps
\(T_{j,\tau_b}f=\tau_{b_j}\cdot f\) meet the approximation
hypothesis by continuity of translations on \(L^2\).
Exact fixed symmetries use \(T_{j,g}f=g\cdot f\).

\subsection{Reflection forms}
For any law of tempered distributions, the linear span of
\(e^{iF(f)}\), \(f\in C_c^\infty(O;V)\) real, is dense in
\(L^2(\Sigma_O)\). To see this, let \(B\) be orthogonal to these
characters and condition it on finitely many evaluations.
Orthogonality to every real linear combination of their tests
says that the Fourier transform of the resulting finite complex
measure is zero. Fourier uniqueness makes that conditional
expectation zero. A countable dense family of compact tests
generates \(\Sigma_O\), since every sample is a continuous
distribution. Conditional expectations on its increasing finite
evaluation sigma-algebras converge in \(L^2\), proving the claim.

Fix \(r\in\mathcal R\). For
\(B(F)=\sum_bz_be^{iF(f_b)}\), the reflected form equals
the left side of \eqref{eq:rp-approx}. Each summand is bounded
and continuous on \(H_{-t}\), so weak convergence gives
\[
 \E_P[\overline{B(r\cdot F)}B(F)]\ge0.
\]
The preceding symmetry argument gives \(r\)-invariance of \(P\).
Consequently
\begin{equation}\label{eq:reflection-form-bound}
 \big|\E_P[\overline{B(r\cdot F)}C(F)]\big|
 \le\|B\|_{L^2(P)}\|C\|_{L^2(P)}.
\end{equation}
Density extends positivity to all
\(B\in L^2(\Sigma_{H_r^+},P)\).
If the plane of \(r\) is \(n\cdot x=0\), its translate at height
\(c\) has reflection
\[
 r_c=\tau_vr\tau_{-v},\qquad v=c|n|^{-2}n.
\]
Translation invariance identifies the two forms and carries their
positive-side sigma-algebras to one another. It also gives
\(r_c\)-invariance.

\begin{proof}[Proof of \cref{thm:structural-limit}]
The exponential bound implies
\(\E_{P_j}F(f)=0\) and
\(\E_{P_j}|F(f)|^2\le\|f\|_2^2\);
see \cref{prop:subgaussian}.
The preceding arguments with \(C=1\) prove the identities,
invariance and reflection positivity.
The Gaussian bound for \(P\) and convergence of all finite-family
moments and cumulants follow from
\cref{thm:gaussian-convergence}, whose hypotheses are
the weak convergence and \eqref{eq:common-gaussian}.
\end{proof}


\section{Tensor Schwinger functions}\label{sec:tensors}
The following estimate is the argument of
\cite[Section~5.3]{QED3Manuscript}. It applies to any
Gaussian-dominated real two-form law in three Euclidean dimensions.
For \(n\ge1\), set
\[
 \mathcal S_n=\mathcal S((\R^3)^n;(\Lambda^2\C^3)^{\otimes n}),\qquad
 q_m(\Psi)=\max_{\substack{|\alpha|\le m\\\text{components}}}
 \sup_y(1+|y|^2)^{m/2}|\partial^\alpha\Psi(y)|.
\]
\begin{proposition}\label{prop:tensor-growth}
Let \(P\) be a law of real tempered two-forms with
\(\E_Pe^{F(f)}\le e^{\|f\|_2^2/2}\) for all real Schwartz \(f\).
For each \(n\ge1\) there is a continuous linear map
\(\Psi\mapsto Q_\Psi\) from \(\mathcal S_n\) to \(L^2(P)\)
which agrees with \(\prod_aF(f_a)\) on decomposable tests.
The functional \(S_n(\Psi)=\E_PQ_\Psi\) is the unique continuous
extension of the smeared moments to \(\mathcal S_n\), and
\begin{equation}\label{eq:tensor-growth}
 |S_n(\Psi)|\le2\,66^n\sqrt{n!}\,q_{6n}(\Psi).
\end{equation}
\end{proposition}

\begin{proof}
Let \(h_\kappa=h_ke_{\ell m}\) be the real Hermite two-form basis.
The eigenvalues of \(A=1+|x|^2-\Delta\) are
\(\lambda_\kappa=2|k|_1+4\). Expand in Schwartz space
\[
 \Psi=\sum_{\boldsymbol\kappa}c_{\boldsymbol\kappa}(\Psi)
 h_{\kappa_1}\otimes\cdots\otimes h_{\kappa_n}.
\]
Counting eigenvalues gives
\[
 Z_4:=\sum_\kappa\lambda_\kappa^{-4}
 =\frac3{32}\sum_{m\ge0}\frac{m+1}{(m+2)^3}
 \le\frac3{32}\left(\frac{\pi^2}{6}-1\right)<\frac1{16}.
\]
Weighted Cauchy--Schwarz and Parseval imply
\[
 \sum_{\boldsymbol\kappa}|c_{\boldsymbol\kappa}(\Psi)|
 \le Z_4^{n/2}\|(A^2)^{\otimes n}\Psi\|_2.
\]
For \(M(x)=1+|x|^2\),
\[
 A^2=M^2-2M\Delta-4x\cdot\nabla-6+\Delta^2.
\]
The sum of the absolute differential coefficients is bounded by
\(34M^2\). Applying this in each argument gives
\[
 |(A^2)^{\otimes n}\Psi(x)|
 \le34^nq_{6n}(\Psi)\prod_{a=1}^n(1+|x_a|^2)^{-1}
\]
componentwise. Since
\(\int_{\R^3}(1+|x|^2)^{-2}\,\dd x=\pi^2\) and there are
\(3^n\) components,
\begin{equation}\label{eq:Hermite-l1}
 \sum_{\boldsymbol\kappa}|c_{\boldsymbol\kappa}(\Psi)|
 \le4^{-n}(34\pi\sqrt3)^nq_{6n}(\Psi).
\end{equation}

For real tests the Gaussian moment estimate is
\[
 \E_P\prod_{a=1}^m|F(f_a)|
 \le c_m\prod_a\|f_a\|_2,\qquad
 c_m=2\,2^{m/2}\Gamma(1+m/2)
 \le2\,2^{m/2}\sqrt{m!};
\]
see \cref{prop:subgaussian} and its proof.
The last inequality is the midpoint log-convexity bound for the
gamma function between \(1\) and \(m+1\).
The \(2n\)-moment estimate on real basis elements shows that
\begin{equation}\label{eq:tensor-random-evaluation}
 Q_\Psi=\sum_{\boldsymbol\kappa}c_{\boldsymbol\kappa}(\Psi)
             \prod_{a=1}^nF(h_{\kappa_a})
\end{equation}
converges absolutely in \(L^2(P)\), with norm at most
\(\sqrt{c_{2n}}\sum|c_{\boldsymbol\kappa}(\Psi)|\).
This proves continuity into \(L^2(P)\).
Expanding each \(f_a\) in the Hermite basis and using the same
product bound identifies the sum with \(\prod_aF(f_a)\) on
decomposable tests. Density of finite Hermite tensors gives
uniqueness of the continuous extension. Finally,
\[
 |S_n(\Psi)|\le c_n\sum_{\boldsymbol\kappa}|c_{\boldsymbol\kappa}(\Psi)|
 \le2\sqrt{n!}\left(\frac{34\pi\sqrt6}{4}\right)^nq_{6n}(\Psi),
\]
which implies \eqref{eq:tensor-growth}.
\end{proof}

\section{Application to the QED3 curvature}\label{sec:qed}
The lattice input is a probability law of real curvature.
For the all-flux chamber model, the infinite-volume construction
in \cite[Theorem~A(i) and Section~2]{ChamberLattice} gives such a
law at each fixed spacing. Its finite model has antiperiodic
spinors, the chamber fermion operator, and parameters
\[
 q>0,\qquad 0<\rho<1,\qquad w_0,w_1,w_2,w_3>0,\qquad M\in\R.
\]
The photon curvature is $u=dA+(2\pi/q)z$, with all closed
integer fluxes $z$ included. The precise operator and photon
measure are those of the cited section. The parameters may
depend on the spacing but are fixed during each volume limit;
the physical charge scaling is $q_a=e\sqrt a$ for fixed $e>0$.

Every selected real-curvature limit $\nu_a$ from that construction
satisfies
\begin{equation}\label{eq:qed-gaussian-input}
 \E_{\nu_a}\exp\left(\sum_{x,i<j}b_{ij}(x)u_{ij}(x)\right)
 \le\exp\left(\frac12\sum_{x,i<j}|b_{ij}(x)|^2\right)
\end{equation}
for every finitely supported real $b$. The constant is one,
independently of the spacing and the selected volume limit.
For completeness, the passage from finite volume uses only a
fixed finite support: it embeds injectively into sufficiently
large tori, so its counting norm is unchanged. The finite Gaussian
estimate \cite[Proposition~5.1]{ChamberLattice} applies there.
Weak convergence passes the bounded continuous cylinder
function $\min\{R,\exp(\sum b_{ij}u_{ij})\}$ to the limit,
and monotone convergence as $R\to\infty$ gives
\eqref{eq:qed-gaussian-input}.

\subsection{Centered plaquettes and their continuum limits}\label{sec:plaquettes}
An oriented lattice two-cochain has real components
\(u_{\ell m}(x)\), \(x\in\Z^3\), \(\ell<m\), with the usual sign
under reversal of orientation. Its coboundary is
\[
 (du)(x)=u_{23}(x+e_1)-u_{23}(x)
        -u_{13}(x+e_2)+u_{13}(x)
        +u_{12}(x+e_3)-u_{12}(x).
\]
For \(p=(x;\ell m)\) and \(a>0\), set
\begin{equation}\label{eq:plaquette-cells}
 Q_a(p)=a\bigl(x+(e_\ell+e_m)/2+[-1/2,1/2)^3\bigr),\qquad
 f_{[a]}(p)=a^{-3/2}\int_{Q_a(p)}f_{\ell m}(y)\,\dd y.
\end{equation}
The \(\ell m\) component of the step field \(F_a\) equals
\(a^{-3/2}u_{\ell m}(x)\) on \(Q_a(p)\).
The geometric action is the exterior-square action on two-forms.


\begin{corollary}\label{cor:qed-continuum}
Let $a_n>0$ tend to zero. For each $n$, let $\nu_{a_n}$ be
a Borel probability law on real arrays
$(u_{ij}(x))_{x\in\Z^3,\,1\le i<j\le3}$ satisfying
\eqref{eq:qed-gaussian-input}. Let \(F_{a_n}\) be their centered-plaquette step fields as in \eqref{eq:plaquette-cells}.
There are a subsequence and a Borel probability law $\mu$ of
real tempered two-forms on $\R^3$ such that the laws of these
step fields converge weakly to $\mu$ in $H_{-4}$. Here
\[
 \|F\|_{-4}^2=\sum_{i<j}\sum_{k\in\N_0^3}
       (2|k|_1+4)^{-8}|F(h_k\,\dd y_i\wedge\dd y_j)|^2.
\]
Each step-field law and every such limit $\mu$ satisfy
\[
 \E e^{F(f)}\le e^{\|f\|_2^2/2}
 \qquad\bigl(f\in\mathcal S(\R^3;\Lambda^2)\text{ real}\bigr).
\]
Along the selected subsequence, all joint moments and cumulants
of complex Schwartz pairings converge to those of $\mu$.
Their finite-family generating functions and logarithms have
the locally uniform convergence asserted in
\cref{thm:gaussian-convergence}.
In particular, this conclusion applies to any choices of the
infinite-volume real-curvature laws of the chamber model just
described.
\end{corollary}

\begin{proof}
For a finitely supported $b$, let $Y=\sum b_{ij}(x)u_{ij}(x)$
and $v^2=\sum|b_{ij}(x)|^2$. Substituting $tb$ in
\eqref{eq:qed-gaussian-input} gives
$\E e^{tY}\le e^{t^2v^2/2}$ for every real $t$.
As in \cite[Proposition~5.1]{QED3Manuscript}, this implies
\begin{equation}\label{eq:qed-second-input}
 \E|Y|^2\le v^2.
\end{equation}
The one-variable calculation is the first paragraph of the
proof of \cref{prop:subgaussian}.

Apply \cref{prop:step-realization} to each law with $d=r=3$,
$C=1$, the three components $12,13,23$, and their offsets
$(e_i+e_j)/2$. Then apply \cref{prop:compactness} to the sequence
of step fields with $s=3$ and $t=4$, as in the proof of
\cref{thm:lattice-continuum}. This also permits different choices
of law when a spacing occurs more than once.
The resulting weak limit on $H_{-4}$ has exactly the
two-form interpretation and norm stated above.

It remains to retain the exponential estimate. For a real
Schwartz two-form $f$, let $b_{a_n}^{(m)}$ be the finite
cell-coefficient truncations used in
\cref{prop:step-realization}. Their pairings converge almost
surely to $F_{a_n}(f)$. Fatou's lemma and
\eqref{eq:qed-gaussian-input} give
\[
 \E e^{F_{a_n}(f)}
 \le\liminf_{m\to\infty}
        \exp\bigl(\|b_{a_n}^{(m)}\|_{\ell^2}^2/2\bigr)
 =\exp\bigl(\|b_{a_n}\|_{\ell^2}^2/2\bigr)
 \le e^{\|f\|_2^2/2}.
\]
This is the infinite-lattice test extension of
\cite[Section~6.1]{LatticeLimits}. Apply
\cref{thm:gaussian-convergence} to the weakly convergent
subsequence already chosen. It gives the Gaussian bound for
$\mu$ and all the stated correlation and cumulant limits.
\end{proof}


The following consequence adds the lattice identities and reflection
hypotheses. The preceding corollary supplies weakly convergent
subsequences; the assertion below applies to every selected weak limit.

\begin{corollary}\label{cor:plaquette-limit}
Let \(a_j\downarrow0\) and let \(\nu_j\) be laws of real
two-cochains on \(\Z^3\). Suppose that:
\begin{enumerate}[label=(\roman*),leftmargin=*]
\item \(du=0\) almost surely and
\(\E_{\nu_j}e^{\langle h,u\rangle}\le e^{\|h\|_{\ell^2}^2/2}\)
for every finitely supported real plaquette test \(h\).
\item Each \(\nu_j\) is invariant under integer translations,
the oriented cubic action and \(u\mapsto-u\).
\item Each \(\nu_j\) is reflection positive for
\(x_\ell=0\) and \(x_\ell\pm x_m=0\), \(\ell<m\),
on character sums supported on plaquettes all of whose corners
lie strictly on the chosen positive side.
\item The laws of \(F_{a_j}\) are realized on \(H_{-4}\) and
converge weakly there to \(P\).
\end{enumerate}
Then \(P\) is a translation-invariant law of closed real tempered
two-forms, invariant under the cubic group and field-sign change.
It satisfies \(\E_Pe^{F(f)}\le e^{\|f\|_2^2/2}\) and is invariant
and reflection positive at every plane
\(x_\ell=c\), \(x_\ell\pm x_m=c\), \(c\in\R\).
All fixed finite-family moments and cumulants converge along
the given sequence.
\end{corollary}

\subsection{Pairings and equivariance}
\Cref{prop:step-realization} gives the measurable realization and
\(F_a(f)=\langle f_{[a]},u\rangle\). The cell estimate
\eqref{eq:cell-contraction} and the Fatou argument in
\cref{cor:qed-continuum} give
\begin{equation}\label{eq:step-gaussian}
 \E e^{F_a(f)}\le e^{\|f\|_2^2/2}.
\end{equation}
If \(\gamma\) is an affine cubic symmetry of the unit lattice
and \(\gamma_a\) its physical rescaling, define
\((\gamma\cdot h)(p)=h(\gamma^{-1}p)\), where reversing the
orientation of a plaquette reverses the coefficient's sign.
A change of variables gives
\begin{equation}\label{eq:limit-test-equivariance}
 (\gamma_a\cdot f)_{[a]}=\gamma\cdot f_{[a]}.
\end{equation}
Signed coordinate permutations carry cell centers to the corresponding
plaquette centers and carry cubes to cubes, up to null boundaries.
The orientation sign is the exterior-square matrix in both actions.
Integer translations give physical translations in \(a\Z^3\).
These are the exact symmetries used in \cref{sec:stability}.

\subsection{The discrete identity and the reflection support}
Let \(\xi\in C_c^\infty(\R^3;\R)\) denote the three-form
\(\xi\,\dd y_1\wedge\dd y_2\wedge\dd y_3\). Define
\[
 \partial_\ell^a\xi(y)
 =a^{-1}\bigl(\xi(y+ae_\ell/2)-\xi(y-ae_\ell/2)\bigr),\qquad
 (v^a_{23},v^a_{13},v^a_{12})
 =(-\partial_1^a\xi,\partial_2^a\xi,-\partial_3^a\xi).
\]
Then
\begin{equation}\label{eq:limit-Bianchi-error}
 F_a(v^a)=0,\qquad
 \|v^a-d^*\xi\|_2
 \le\frac{a^2}{24}\sum_{\ell=1}^3\|\partial_\ell^3\xi\|_2.
\end{equation}
For the exact identity, assign to each unit cube the number
\[
 \xi_{[a]}(x)=a^{-5/2}\int_{a(x+[0,1)^3)}\xi(y)\,\dd y.
\]
The \(23,13,12\) coefficients of \(v^a_{[a]}\) are
\[
 -(\xi_{[a]}(x)-\xi_{[a]}(x-e_1)),\quad
 \xi_{[a]}(x)-\xi_{[a]}(x-e_2),\quad
 -(\xi_{[a]}(x)-\xi_{[a]}(x-e_3)).
\]
Their finite pairing with \(u\) regroups to
\(\sum_x\xi_{[a]}(x)(du)(x)=0\).
The Fourier multiplier of
\(\partial_\ell^a-\partial_\ell\) is
\(i(2a^{-1}\sin(ap_\ell/2)-p_\ell)\), whose modulus is at most
\(a^2|p_\ell|^3/24\). Plancherel gives the second assertion.

For a mirror with defining function
\(\ell(y)=y_j\) or \(y_j\pm y_k\), take compact tests supported
in \(\{\ell\ge\delta\}\), \(\delta>0\).
A nonzero cell coefficient comes from a cell meeting this set.
For every corner \(q\) of its physical plaquette and every \(y\)
in that cell, \(|\ell(q)-\ell(y)|\le2a\):
each coordinate displacement is at most \(a\), and the normal
has at most two nonzero coordinates of modulus one.
Thus \(2a<\delta\) puts all contributing plaquette corners
strictly on the positive side. There are only finitely many such
plaquettes. The lattice reflection inequality and
\eqref{eq:limit-test-equivariance} give \eqref{eq:rp-approx}.

\begin{proof}[Proof of \cref{cor:plaquette-limit}]
Equation \eqref{eq:step-gaussian} gives the common Gaussian input.
Equation \eqref{eq:limit-Bianchi-error} verifies the identity
hypothesis with \(D^*=d^*\). Equivariance gives cubic symmetry and
mesh translations converging to every continuum translation.
The support check gives the eventual character inequalities at the
nine origin mirrors. \Cref{thm:structural-limit} applies.
Field-sign invariance passes directly through its
characteristic-function identity.
\end{proof}

The real-curvature laws supplied by \cite[Theorem~A]{ChamberLattice}
satisfy the additional hypotheses of \cref{cor:plaquette-limit}.
Together with \cref{cor:qed-continuum}, this gives subsequential
continuum laws with all the stated structural properties. The
reflection passage is the origin-height argument of
\cite[Proposition~6.1]{QED3Manuscript}; translations supply every
height. The conditional corollary imposes no arithmetic restriction
on the spacing sequence.

\subsection{Recovery of a prescribed volume limit}
\label{subsec:fixed-spacing-recovery}\label{sec:recovery}
An infinite-lattice law gives a step-field law by
\cref{prop:step-realization}. To treat a prescribed volume limit
of the finite-torus fields, we also need the converse
identification. We use the bump-function argument of
\cite[Section~6.4]{QED3Manuscript}.

Fix a spacing $a>0$. Let $\mathcal P_2$ be the positively oriented
plaquettes of $\Z^3$, and put
\[
 \Omega=\{u\in\R^{\mathcal P_2}:du=0\},
\]
with its coordinate product topology. For each integer $n\ge4$,
let $\lambda_{a,n}$ be a probability law of real closed
two-cochains on $(\Z/n\Z)^3$ satisfying
\begin{equation}\label{eq:recovery-finite-gaussian}
 \E_{\lambda_{a,n}}e^{\langle\beta,u\rangle}
 \le e^{\|\beta\|_{\ell^2((\Z/n\Z)^3)}^2/2}
 \qquad(\beta\text{ a real torus two-cochain}).
\end{equation}
Both the pairing and the norm sum over positively oriented torus
plaquettes. Denote the law of the periodic extension by
$\nu_{a,n}$. Use the centered interpolation and coefficients
$f_{[a]}$ of \eqref{eq:plaquette-cells}, and write $P_{a,n}$
for its law on $H_{-4}$.
For general Schwartz tests these finite-torus laws obey the
periodized bound \eqref{eq:limit-cutoff-G} below, with physical
period $R=an$. The following identification gives the sharp
$L^2$ bound after the volume limit.

\begin{proposition}\label{prop:fixed-spacing-recovery}
Suppose $n_j\to\infty$ and $P_{a,n_j}\Longrightarrow\mu_a$
weakly on $H_{-4}$. Then there is a probability law $\nu_a$
on $\Omega$ such that
\begin{equation}\label{eq:limit-volume-identification}
 \nu_{a,n_j}\Longrightarrow\nu_a,\qquad
 \mu_a=(\sharp_a)_*\nu_a,
\end{equation}
where $\sharp_a$ is the centered step map, defined as a Borel
map into $H_{-4}$ on a set of full $\nu_a$ measure and set
equal to zero elsewhere. The convergence of cochain laws is
in the product topology and holds along the given sequence.
For every finitely supported real cochain $\beta$,
\begin{equation}\label{eq:limit-lattice-G}
 \E_{\nu_a}e^{\langle\beta,u\rangle}
 \le e^{\|\beta\|_{\ell^2(\mathcal P_2)}^2/2}.
\end{equation}
In particular, $\mu_a$ satisfies the Gaussian bound
$\E_{\mu_a}e^{F(f)}\le e^{\|f\|_2^2/2}$ for real Schwartz
two-forms $f$.
\end{proposition}

\begin{proof}
Discard finitely many terms so that $an_j\ge1$.
The finite-volume step maps take values in $H_{-3}$.
Indeed, the periodization estimate below and the Hermite bound
$p_1(h_ke_\alpha)\le31\lambda_k$ from
\cite[Section~3]{QED3Integration} give
$|F_a(h_ke_\alpha)|\le31\lambda_k\|u\|_{\ell^2}$.
Summing with weight $\lambda_k^{-6}$ proves continuity of
each finite-dimensional step map into $H_{-3}$, since
$\sum_k\lambda_k^{-4}<\infty$.

We use the realization already proved in
\cref{prop:step-realization}. Let $\Omega_3$ be its weighted
square-summability domain with exponent $p=3$, intersected with
$\Omega$. For every $u\in\Omega_3$, the step distribution has the
absolutely convergent pairings \eqref{eq:absolute-pairing}.
Let $\Omega^{(a)}\subset\Omega_3$ be the set on which it has
finite $H_{-3}$ norm, and set $\sharp_a=0$ off this set.
The measurability argument in \cref{prop:compactness} makes
$\Omega^{(a)}$ measurable and $\sharp_a$ a Borel map into $H_{-4}$.
Every periodic cochain lies in $\Omega_3$ and, by the preceding
finite-volume estimate, in $\Omega^{(a)}$.

For each plaquette $\sigma$, choose a real smooth bump $b^\sigma$
with only its corresponding two-form component nonzero, supported
inside $Q_a(\sigma)$ and of integral $a^{3/2}$. Then
$(b^\sigma)_{[a]}=\mathbf1_\sigma$. The map
\begin{equation}\label{eq:limit-step-inverse}
 \mathcal J_a:H_{-4}\longrightarrow\R^{\mathcal P_2},
 \qquad \mathcal J_a(F)(\sigma)=F(b^\sigma),
\end{equation}
is continuous for the product topology, and
$\mathcal J_a\sharp_a u=u$ on $\Omega^{(a)}$.
At finite volume the identities are
\[
 P_{a,n_j}=(\sharp_a)_*\nu_{a,n_j},\qquad
 \nu_{a,n_j}=(\mathcal J_a)_*P_{a,n_j}.
\]
The mapping theorem therefore gives
$\nu_{a,n_j}\Longrightarrow\nu_a:=(\mathcal J_a)_*\mu_a$.
The closed conditions $du=0$ pass to this product limit.

Fix a finitely supported real $\beta$ on the covering lattice.
For all sufficiently large $n_j$, its support embeds injectively
in the torus. Its folded norm then equals its covering
$\ell^2$ norm. Apply \eqref{eq:recovery-finite-gaussian}, pass
the bounded continuous cylinder
$\min\{R,e^{\langle\beta,u\rangle}\}$ to the limit, and let
$R\to\infty$. This proves \eqref{eq:limit-lattice-G}.

The Gaussian bound implies the finite-coefficient second-moment
bound, so \cref{prop:step-realization} gives $\nu_a(\Omega_3)=1$.
The step distribution on this domain satisfies the common
$L^2$ test estimate. Applying \cref{prop:compactness} with $s=3$
therefore gives $\nu_a(\Omega^{(a)})=1$.
The finite-coefficient Fatou argument in
\cref{cor:qed-continuum} now gives
\[
 \E_{\nu_a}e^{(\sharp_a u)(f)}
 \le e^{\|f_{[a]}\|_{\ell^2}^2/2}
 \le e^{\|f\|_2^2/2}.
\]

It remains to prove the reverse equality in
\eqref{eq:limit-volume-identification}. Choose a countable
Schwartz-dense family $\mathcal D_0$ of real compact smooth
two-forms. For $f\in\mathcal D_0$, the test
\[
 f'=f-\sum_\sigma f_{[a]}(\sigma)b^\sigma
\]
is a finite sum and satisfies $f'_{[a]}=0$. Every finite-volume
step field therefore annihilates $f'$. The hyperplane
$\{F:F(f')=0\}$ is closed in $H_{-4}$ and has full $\mu_a$
measure. Intersect these events for $f\in\mathcal D_0$ with
the full-measure event
$\mathcal J_a(F)\in\Omega^{(a)}$. On that single event,
\[
 F(f)=\sum_\sigma f_{[a]}(\sigma)F(b^\sigma)
      =(\sharp_a\mathcal J_aF)(f)
 \qquad(f\in\mathcal D_0).
\]
Both sides are tempered distributions, so density gives
$F=\sharp_a\mathcal J_aF$ there. Taking laws proves the
pushforward identity and the Gaussian assertion for $\mu_a$.
\end{proof}

\subsection{Finite-torus estimates and joint cutoff sequences}
\label{subsec:joint-cutoff-sequences}
The exact $L^2$ Gaussian bound in
\cref{thm:gaussian-convergence} applies to the recovered
infinite-volume laws. Before the volume limit, the Schwartz
test is periodized. We recall the estimate from
\cite[Sections~6.1--6.2]{QED3Manuscript} that controls this step.

Put $R=an$ and $f^R(y)=\sum_{k\in\Z^3}f(y+Rk)$.
The finite-torus coefficient is
\begin{equation}\label{eq:folded-cell-coefficient}
 b_{a,n}([\sigma])
 =\sum_{k\in\Z^3}f_{[a]}(\sigma+nk)
 =a^{-3/2}\int_{Q_a(\sigma)}f^R_{\ell m}(y)\,\dd y
 \quad(\sigma=(x;\ell m)).
\end{equation}
The sum converges absolutely. It gives
$F_a(f)=\langle b_{a,n},u\rangle$, and cell Cauchy--Schwarz
gives $\|b_{a,n}\|_{\ell^2}\le
\|f^R\|_{L^2(\R^3/R\Z^3)}$. With $Q_R=[-R/2,R/2)^3$, define
\begin{equation}\label{eq:limit-periodic-seminorm}
 p_R(f)=\|f\|_2+
        26R^{-2}\||y|^2f\|_{L^2(\R^3\setminus Q_R)}.
\end{equation}
Then
\begin{equation}\label{eq:limit-cutoff-G}
 \E_{P_{a,n}}e^{F(f)}
 \le e^{\|b_{a,n}\|_{\ell^2}^2/2}
 \le e^{p_R(f)^2/2}
 \qquad(f\text{ real}).
\end{equation}
To see the last estimate, put
$c_k=\|f\|_{L^2(Q_R+Rk)}$. Minkowski gives
$\|f^R\|_{L^2(\R^3/R\Z^3)}\le\sum_kc_k$. For $k\ne0$,
$|y|\ge R|k|_\infty/2$ on $Q_R+Rk$, and weighted
Cauchy--Schwarz gives
\[
 \sum_{k\ne0}c_k
 \le4R^{-2}
   \left(\sum_{k\ne0}|k|_\infty^{-4}\right)^{1/2}
   \||y|^2f\|_{L^2(\R^3\setminus Q_R)}.
\]
The shell count $24j^2+2$ gives
$\sum_{k\ne0}|k|_\infty^{-4}=4\pi^2+\pi^4/45$,
whose square root times four is less than $26$.
For $R\ge1$, one has $p_R(f)\le p_1(f)$, and
$p_R(f)\to\|f\|_2$ as $R\to\infty$. The seminorm $p_1$
is continuous on Schwartz space.

\begin{corollary}\label{cor:prescribed-volume-first}
At each spacing $a_i\downarrow0$, take the chamber torus laws
with parameters fixed during that volume limit. Let $\mu_i$
be any prescribed weak $H_{-4}$ volume limit of their
centered step fields. Suppose $\mu_i\Longrightarrow\mu$ on
$H_{-4}$. Then $\mu$ has the conclusions of
\cref{cor:plaquette-limit}, including convergence of all
finite-family moments and cumulants from $\mu_i$ to $\mu$.
Its tensor Schwinger functions satisfy \cref{prop:tensor-growth}.
There are torus sides $n_i\to\infty$ such that
$a_in_i\to\infty$,
\[
 P_{a_i,n_i}\Longrightarrow\mu,
\]
so $\mu$ is also realized by a joint sequence of finite
cutoff laws.
\end{corollary}

\begin{proof}
Apply \cref{prop:fixed-spacing-recovery} at each $a_i$.
It identifies the prescribed law $\mu_i$ with the step image
of a weak cochain limit. The chamber theorem
\cite[Theorem~A]{ChamberLattice} applies to every such limit
with the parameters fixed: it gives closedness, Gaussian
domination, the lattice symmetries, and reflection positivity
at the nine mirror families. Thus these recovered laws meet
all the hypotheses of \cref{cor:plaquette-limit}; that
corollary and \cref{prop:tensor-growth} give the assertions
about $\mu$ and the sequence $\mu_i$.

Fix a metric for weak convergence of probability laws on
$H_{-4}$. Along the defining volume sequence for $\mu_i$,
choose $n_i$ so large that $a_in_i\ge i$ and the distance
between $P_{a_i,n_i}$ and $\mu_i$ is less than $1/i$.
Then $P_{a_i,n_i}\Longrightarrow\mu$ and
$R_i:=a_in_i\to\infty$. This is the diagonal selection in
\cite[proof of Proposition~6.1]{QED3Manuscript}.
\end{proof}



\printbibliography
\end{document}
